\documentclass[11pt]{article}
\usepackage[T1]{fontenc}
\usepackage{lmodern}
\usepackage[utf8]{inputenc}
\usepackage{amsmath,amssymb,amsthm,mathtools}
\usepackage{microtype}
\usepackage[margin=1in]{geometry}
\usepackage{booktabs}
\usepackage{array}
\usepackage{enumitem}
\usepackage{xcolor}
\usepackage{listings}
\usepackage{hyperref}
\usepackage{fancyhdr}
\usepackage{lastpage}
\usepackage{longtable}
\usepackage{caption}

\hypersetup{
  colorlinks=true,
  linkcolor=blue!45!black,
  urlcolor=blue!45!black,
  citecolor=blue!45!black,
  pdftitle={Contender 14: Self-Indexed Diagonal Operator Reflection},
  pdfauthor={Research proposal for BusyLean}
}

\pagestyle{fancy}
\fancyhf{}
\lhead{BusyLean Contender 14}
\rhead{Diagonal operator reflection}
\cfoot{\thepage\ of \pageref{LastPage}}
\setlength{\headheight}{14pt}

\lstdefinelanguage{lean}{
  keywords={def,theorem,inductive,structure,where,match,with,fun,let,in,
    namespace,end,open,abbrev,import,by,have,show,exact,simpa,using,
    induction,termination_by,if,then,else},
  sensitive=true,
  morecomment=[l]{--},
  morecomment=[s]{/-}{-/},
  morestring=[b]"
}
\lstset{
  language=lean,
  basicstyle=\ttfamily\small,
  keywordstyle=\bfseries\color{blue!45!black},
  commentstyle=\itshape\color{black!60},
  stringstyle=\color{green!35!black},
  columns=fullflexible,
  keepspaces=true,
  frame=single,
  framerule=0.3pt,
  rulecolor=\color{black!25},
  xleftmargin=1em,
  xrightmargin=1em,
  breaklines=true,
  breakatwhitespace=false,
  showstringspaces=false,
  literate={→}{{$\to$}}1 {←}{{$\leftarrow$}}1 {ℕ}{{$\mathbb{N}$}}1
    {∀}{{$\forall$}}1 {≤}{{$\le$}}1 {≥}{{$\ge$}}1 {×}{{$\times$}}1
    {⟨}{{$\langle$}}1 {⟩}{{$\rangle$}}1 {λ}{{$\lambda$}}1
    {∃}{{$\exists$}}1 {∈}{{$\in$}}1 {ω}{{$\omega$}}1
}

\newtheorem{theorem}{Theorem}[section]
\newtheorem{lemma}[theorem]{Lemma}
\newtheorem{corollary}[theorem]{Corollary}
\theoremstyle{definition}
\newtheorem{definition}[theorem]{Definition}
\newtheorem{construction}[theorem]{Construction}
\theoremstyle{remark}
\newtheorem{remark}[theorem]{Remark}
\newtheorem{warning}[theorem]{Warning}

\newcommand{\N}{\mathbb{N}}
\newcommand{\lean}[1]{\texttt{#1}}
\newcommand{\Fam}{\mathsf{Fam}}
\newcommand{\Trans}{\mathsf{Transformer}}
\newcommand{\Prog}{\mathsf{Progressive}}
\newcommand{\Fd}{\mathsf{F}_{\mathsf{dyn}}}
\newcommand{\opstage}{\mathsf{OpStage}}
\newcommand{\tower}{\mathsf{Tower}}
\newcommand{\sched}{\mathsf{Schedule}}
\newcommand{\C}[1]{C_{#1}}
\newcommand{\F}[1]{F_{#1}}
\newcommand{\Rfam}{R_{\Fd}^{\mathrm{fam}}}
\newcommand{\code}[1]{\texttt{#1}}

\title{\textbf{Contender 14: Self-Indexed Diagonal Operator Reflection}\\[0.35em]
\large A constructive large-number proposal for BusyLean}
\author{Research proposal prepared for Vladimir Reshetnikov}
\date{October 13, 2026}

\begin{document}
\maketitle

\begin{abstract}
We propose a new closed computable natural number for the BusyLean contest,
targeting repository commit
\texttt{1ee50c7d97bbed57a60709f661453c3e213f4db3}, whose incumbent is
Contender~13.  Contender~13 performs one reflection step over a predicative
``dynamic-kind'' extension of System~$F_\omega$, with Contender~12's family as
the unary oracle.  The proposed Contender~14 makes three independent upgrades.
First, it lifts dynamic-kind reflection from a unary oracle to a packed
countable oracle family, while retaining an explicit prefix maximum that gives
direct strict inequalities without any false monotonicity assumption.  Second,
it turns this family reflector into a progressive transformer and iterates the
repository's operator-stage construction through rank-dependent autonomous
schedules.  Third, it diagonalizes the rank: at input $n$, the number of
operator-stage levels is
\[
  \mathsf{family13}(n)=F_{13}(F_{13}(n)),
\]
so at input $42$ the tower rank is definitionally the entire previous
contender.  The executable definition nevertheless does not mention the
constant \lean{contender\_13}; it uses the predecessor's construction, as the
contest rules permit.

The supplied Lean development states and proves, at source level, that every
finite tower rank is progressive, every lower rank is strictly below every
higher rank along explicit embedded branches, every enumerated dynamic-kind
program at depth $n$ is below $F_{14}(n)$, $\mathsf{family13}(n)<F_{14}(n)$
for every $n$, and finally
\[
  \mathsf{contender\_13}<\mathsf{contender\_14}.
\]
All new definitions are transparent and structurally recursive, and the
archive includes an integration patch, static audit, and reproduction notes.
The current execution environment did not contain Lean~4.34.0 or the repository
checkout, so kernel compilation and timing remain mandatory integration checks;
this limitation is recorded explicitly rather than hidden.
\end{abstract}

\tableofcontents

\section{Executive statement of the proposal}

The proposal consists of the following definitions.

\begin{align*}
\Rfam(fs,n)
  &= n+1+\max\!\left(
       \max_{i\le n}fs(i,n),
       \max_{t\in\mathrm{CTm}_{\le n}}
          \mathrm{value}_{\Fd}(\mathrm{pack}(fs),n,t)
     \right),\\
T_0&=\text{the transformer induced by }\Rfam,\\
T_{k+1}&=\opstage(T_k,\sched(k)),\\
\mathsf{family13}(n)&=F_{13}(F_{13}(n)),\\
F_{14}(n)&=T_{\mathsf{family13}(n)}
          (\text{constant family }\mathsf{family13})(n,n),\\
\C{14}&=F_{14}(F_{14}(42)).
\end{align*}

The rank-dependent schedule $\sched(k)$ has one deliberately simple branch and
countably many very large branches:
\[
\sched(k)=\operatorname{node}\!\left(j\mapsto
  \begin{cases}
  \operatorname{leaf},&j=0,\\
  \operatorname{realize}\!\left(
    \operatorname{hyper}_{k+r+1}(\operatorname{encode}(\mathsf{tree12}))
  \right),&j=r+1.
  \end{cases}
\right).
\]
The leaf is not cosmetic.  It is a proof witness: it embeds one ordinary
autonomous lift inside each new operator stage, allowing a direct strict
comparison even though reflection is not monotone in its oracle.  The other
branches make the schedule genuinely diagonal and unbounded in the explicit
higher-derivative notation already present in the repository.

The final Lean line is intentionally simple:

\begin{lstlisting}
def contender_14 : Nat := F14 (F14 42)
\end{lstlisting}

The size comes from the construction of \lean{F14}, not from an enormous
literal or an unproved external oracle.

\section{Repository baseline and contest constraints}

\subsection{Baseline}

This work targets the BusyLean main branch at the following commit:
\begin{center}
\texttt{\small 1ee50c7d97bbed57a60709f661453c3e213f4db3}.
\end{center}
At that commit the repository uses Lean~4.34.0 and declares Contender~13 as the
current entry:

\begin{lstlisting}
def family12 (n : Nat) : Nat := F12 (F12 n)
def F13 (n : Nat) : Nat := Fd.reflectFd family12 n
def contender_13 : Nat := F13 (F13 42)
\end{lstlisting}

The repository already contains four ingredients needed here:

\begin{enumerate}[leftmargin=*]
\item \lean{Fd.reflectFd}, a total reflector for the predicative dynamic-kind
  language;
\item \lean{packFamily}, an executable coding of a countable family of unary
  functions as one unary oracle;
\item \lean{operatorStage}, a generic closure operation on family
  transformers;
\item \lean{Autonomous.hyper}, \lean{encode}, and \lean{realize}, which build
  transparent countably branching schedules above an existing tree.
\end{enumerate}

The proposal composes these existing certified mechanisms rather than adding a
new unverified interpreter or a new normalization theorem.

\subsection{Relevant rules}

The construction is designed around the repository's stated rules:

\begin{itemize}[leftmargin=*]
\item the number is a closed, constructive \lean{Nat};
\item the previous contender must be beaten by a formal theorem;
\item no \lean{noncomputable}, \lean{partial}, \lean{unsafe}, custom axiom,
  \lean{opaque} stand-in, \lean{extern}, \lean{implemented\_by},
  \lean{native\_decide}, or \lean{sorry} is used;
\item only the inherited permitted axioms \lean{propext} and
  \lean{Quot.sound} may occur in proof closures;
\item the previous contender constant must not be used as a cheap numerical
  parameter, although its underlying family construction may be used as an
  oracle, seed, or component of a qualitatively stronger mechanism;
\item definitions and proofs must satisfy the repository's 15-second and
  60-second phase budgets on the designated toolchain.
\end{itemize}

The last item is an empirical requirement.  Source inspection can establish
structural computability and dependency discipline, but only an actual clean
Lean build can establish elaboration, kernel acceptance, axiom reports, and
wall-clock timing.

\section{Why a new layer is available above Contender 13}

Contender~13 is a major language-strength jump.  Its object language contains a
predicative dynamic kind interpreted semantically as a dependent sum of
operators at all dynamic-free kinds.  It can quantify over two levels of kind
variables and climb finite extensions above the dynamic kind at run time.
Nevertheless, the numerical wrapper around this language is still a direct
unary reflection:
\[
F_{13}(n)=\operatorname{reflectFd}(\mathsf{family12},n).
\]
This leaves two orthogonal amplification directions.

\paragraph{Family access.}
A unary oracle gives one function $\N\to\N$.  A forest or operator stage needs a
countable family $fs:\N\to\N\to\N$ so an object-language program can select a
child or coordinate dynamically.  The earlier Contenders~9 and~11 solve this
by packing $(i,m)$ into one natural input.  Applying the same idea to the
stronger dynamic-kind language yields a family-aware dynamic reflector.

\paragraph{Operator closure.}
Contender~12 already showed that a progressive reflector can be lifted from a
single family transformation to an operator-stage closure along large
schedules.  Contender~13 strengthened the object language but returned to one
direct reflection step.  The most economical next move is therefore not a new
syntax of doubtful consistency, but to feed the stronger reflector into the
existing generic operator machinery and then diagonalize the number of such
closures.

The construction is consequently modular:
\[
\begin{gathered}
\boxed{\text{stronger object language}}
\quad+\quad
\boxed{\text{family reflection}}\\[0.5em]
\boxed{\text{operator-stage tower}}
\quad+\quad
\boxed{\text{self-indexed rank}}.
\end{gathered}
\]

\section{Family-aware dynamic-kind reflection}

\subsection{Definition}

Let \(\Fam=\N\to\N\to\N\).  The repository's \lean{packFamily} maps
$fs\in\Fam$ to a unary oracle by decoding an input into a pair $(i,m)$ and
returning $fs(i,m)$.  We define

\begin{definition}[Family-aware dynamic reflector]
For $fs\in\Fam$ and $n\in\N$, let
\begin{align*}
P(fs,n)&=\max_{i\le n}fs(i,n),\\
Q(fs,n)&=\max_{t\in\mathrm{CTm}_{\le n}}
  \mathrm{fvalue}_{\Fd}(\mathrm{packFamily}(fs),n,t),\\
\Rfam(fs,n)&=n+1+\max(P(fs,n),Q(fs,n)).
\end{align*}
\end{definition}

The exact Lean definition is:

\begin{lstlisting}
def reflect14 (fs : W.Family) (n : Nat) : Nat :=
  n + 1 + BigNumber.maxN (W.prefixMax (fun i => fs i n) n)
    (BigNumber.maximum (Fd.fvalue (packFamily fs) n) (Fd.ctmsUpTo n))
\end{lstlisting}

There are two logically distinct maxima.

\begin{itemize}[leftmargin=*]
\item The prefix maximum provides a direct, syntax-independent inequality for
  every visible family coordinate.
\item The program maximum lets every accepted dynamic-kind program of bounded
  syntax/fuel query the packed family oracle and be exceeded.
\end{itemize}

Neither is redundant.  Packing gives expressive access to the whole family,
but \lean{Fd.reflectFd} by itself only visibly exceeds the unary oracle at its
current input.  Conversely, the prefix maximum alone would dominate coordinates
but would not expose the family to object-language computation.

\subsection{Direct inequalities}

The arithmetic proof follows the pattern already used by the repository's
other reflectors.

\begin{lemma}[Arithmetic reflection bound]
If
\[
m\le\max(P(fs,n),Q(fs,n)),
\]
then $m<\Rfam(fs,n)$.
\end{lemma}

\begin{proof}
The defining right-hand side is $n+1+M$ with $M=\max(P,Q)$.  From $m\le M$ and
$M<M+1\le n+1+M$, transitivity gives $m<\Rfam(fs,n)$.
\end{proof}

\begin{corollary}[Visible-family domination]
For every $i\le n$,
\[
fs(i,n)<\Rfam(fs,n).
\]
\end{corollary}

\begin{proof}
The finite prefix maximum contains $fs(i,n)$, so
$fs(i,n)\le P(fs,n)\le\max(P,Q)$; apply the arithmetic bound.
\end{proof}

\begin{corollary}[Program domination]
For every $t\in\mathrm{CTm}_{\le n}$,
\[
\mathrm{fvalue}_{\Fd}(\mathrm{packFamily}(fs),n,t)<\Rfam(fs,n).
\]
\end{corollary}

\begin{proof}
The finite maximum $Q(fs,n)$ contains the program value; apply the arithmetic
bound.
\end{proof}

\subsection{Why no monotonicity is used}

A recurring trap in large-number reflection arguments is the claim that making
an oracle pointwise larger can only make every program value larger.  That is
false for ordinary programmable oracles: a program may compare, branch, or
return a small result when an oracle value crosses a threshold.  The BusyLean
repository correctly avoids this assumption, and Contender~14 does the same.
Every strict comparison follows an explicit occurrence in a finite maximum or
an explicit branch in a tree.  This discipline is essential because the later
operator construction nests oracle-dependent interpreters many times.

\section{The base transformer}

A transformer maps one oracle family to another:
\[
\Trans=\Fam\to\Fam.
\]
The family-aware reflector becomes
\[
T_0(fs)(o,n)=\Rfam(fs,n),
\]
where the output coordinate $o$ is deliberately ignored.

\begin{lstlisting}
def fdFamilyTransformer : Transformer := fun fs _ n => reflect14 fs n
\end{lstlisting}

The generic operator library uses the following invariant.

\begin{definition}[Progressive transformer]
A transformer $T$ is progressive when
\[
\forall fs,i,o,n,
\quad i\le n\Longrightarrow fs(i,n)<T(fs)(o,n).
\]
\end{definition}

\begin{theorem}[Base progressivity]
\lean{fdFamilyTransformer} is progressive.
\end{theorem}

\begin{proof}
For visible $i\le n$, the desired inequality is exactly visible-family
domination.  The arbitrary output coordinate disappears because the base
transformer ignores it.
\end{proof}

This one-line invariant is what makes the rest of the construction reusable.
No semantic property of the dynamic-kind language is needed beyond the already
certified total evaluator used inside \lean{Fd.fvalue}.

\section{Rank-dependent diagonal schedules}

\subsection{Existing autonomous notation machinery}

The repository contains a transparent tree syntax with zero, successor, and
countable limit nodes, together with operations suggestive of ordinal addition,
multiplication, $\omega$-powers, derivatives, a finite critical hierarchy, a
Gamma-like diagonal closure, and higher derivatives.  It deliberately treats
these as tree syntax rather than asserting quotient-level ordinal identities.
The functions \lean{encode} and \lean{realize} convert between this syntax and
the W-tree schedules used by reflection.

Contender~12 defines \lean{tree12} by applying 42 higher derivatives above the
language-defined index.  We use \lean{tree12} as a baseline and then build a
schedule family rather than one fixed tree.

\subsection{The schedule}

For rank $k$, define
\[
\sched(k)=\operatorname{node}(b_k)
\]
with
\[
b_k(j)=
\begin{cases}
\operatorname{leaf},&j=0,\\
\operatorname{realize}\bigl(
  \operatorname{hyper}_{k+r+1}(\operatorname{encode}(\mathsf{tree12}))
\bigr),&j=r+1.
\end{cases}
\]

The Lean code is:

\begin{lstlisting}
def schedule14 (k : Nat) : WTree Nat :=
  .node (fun j =>
    match j with
    | 0 => .leaf
    | r + 1 =>
        Autonomous.realize
          (Autonomous.hyper (k + r + 1) (Autonomous.encode tree12)))
\end{lstlisting}

The two branch classes have different jobs.

\paragraph{Branch zero: a formal witness.}
The leaf ensures that the operator stage based on \lean{schedule14 k}
contains an explicit visible leaf-stage child.  By definition, this child is
\lean{liftTransformer base}.  The generic theorem
\lean{transformer\_lt\_lift\_at} then gives a strict path
\[
T(fs)(i,n)
 < \operatorname{lift}(T)(fs)(0,n)
 = \operatorname{OpStage}(T,\operatorname{leaf})(fs)(0,n)
 < \operatorname{OpStage}(T,\sched(k))(fs)(o,n).
\]
This is a direct structural embedding, not a monotonicity argument.

\paragraph{Positive branches: diagonal growth.}
Branch $r+1$ uses a higher-derivative level $k+r+1$.  Thus each fixed rank $k$
contains countably many schedules of unbounded explicit hyper-level, and the
minimum positive level also increases with $k$.  No exact proof-theoretic
ordinal is claimed.  The formal role of these branches is to make the schedule
genuinely richer; the strict inequality proof needs only branch zero.

\section{The diagonal operator tower}

\subsection{Definition}

The generic \lean{operatorStage} accepts a base transformer and a W-tree and
returns a new transformer.  We iterate it with a schedule that changes at every
rank:

\begin{lstlisting}
def diagonalOperatorTower (schedule : Nat → WTree Nat)
    (base : Transformer) : Nat → Transformer
  | 0 => base
  | k + 1 =>
      operatorStage (diagonalOperatorTower schedule base k) (schedule k)
\end{lstlisting}

In mathematical notation,
\[
T_0=T_{\mathrm{base}},\qquad
T_{k+1}=\opstage(T_k,\sched(k)).
\]
This is a genuine iteration of an operator closure: the entire transformer at
rank $k$ becomes the base transformer at rank $k+1$.

\subsection{Progressivity of every rank}

\begin{theorem}[Tower progressivity]
If $T_0$ is progressive, then every $T_k$ is progressive.
\end{theorem}

\begin{proof}
Induct on $k$.  The zero case is the hypothesis.  At a successor, the
repository theorem \lean{operatorStage\_progressive} says that applying an
operator stage to any progressive base transformer yields a progressive
transformer, for every schedule tree.  Instantiate it with $T_k$ and
$\sched(k)$.
\end{proof}

The Lean proof is correspondingly short:

\begin{lstlisting}
theorem diagonalOperatorTower_progressive
    (schedule : Nat → WTree Nat) (base : Transformer)
    (hbase : Progressive base) (k : Nat) :
    Progressive (diagonalOperatorTower schedule base k) := by
  induction k with
  | zero => simpa [diagonalOperatorTower] using hbase
  | succ k ih =>
      simpa [diagonalOperatorTower] using
        (operatorStage_progressive
          (diagonalOperatorTower schedule base k) ih (schedule k))
\end{lstlisting}

\subsection{Every successor rank is strictly larger}

Progressivity alone says that each rank exceeds its seed family.  The canonical
leaf branch proves more: rank $k+1$ directly and strictly embeds rank $k$.

\begin{theorem}[Strict rank step]
For every seed family, every $k$, every $i\le n$, and every output coordinate
$o$,
\[
T_k(\mathrm{seed})(i,n)<T_{k+1}(\mathrm{seed})(o,n).
\]
\end{theorem}

\begin{proof}
By tower progressivity, $T_k$ is progressive.  The explicit path through
branch zero of $\sched(k)$ gives
\[
T_k(\mathrm{seed})(i,n)
 < \operatorname{lift}(T_k)(\mathrm{seed})(0,n)
 < \opstage(T_k,\sched(k))(\mathrm{seed})(o,n).
\]
The rightmost term is definitionally $T_{k+1}(\mathrm{seed})(o,n)$.
\end{proof}

The successor theorem admits a finite transitive closure.

\begin{theorem}[Strict comparison of arbitrary ranks]
For every seed family, $a<b$, $i\le n$, and output coordinate $o$,
\[
T_a(\mathrm{seed})(i,n)<T_b(\mathrm{seed})(o,n).
\]
\end{theorem}

\begin{proof}
Induct on $b$.  The zero case is impossible.  For $b=c+1$, either $a=c$, in
which case the successor theorem applies directly, or $a<c$.  In the latter
case, use the induction hypothesis to reach coordinate zero of rank $c$, then
use one successor step from rank $c$ to rank $c+1$.  Every intermediate
coordinate is explicitly visible, so no monotonicity premise is introduced.
\end{proof}

This is the formal large-margin certificate.  The final construction does not
merely wrap Contender~13 once: its rank parameter at input 42 is the previous
contender, and the source proves a strict comparison between every pair of
ranks below that height.

\section{Self-indexing by the previous construction}

\subsection{The predecessor family}

Following the repository's convention, expose the full Contender~13
construction as a function of its input:

\begin{lstlisting}
def family13 (n : Nat) : Nat := F13 (F13 n)
\end{lstlisting}

At input 42 this is definitionally the previous contender:
\[
\mathsf{family13}(42)=F_{13}(F_{13}(42))=\C{13}.
\]
The new executable definition refers to \lean{F13} and \lean{family13}, not to
the closed constant \lean{contender\_13}.  This matters for the no-laziness
audit: the prior construction is used as a seed and meta-rank for a stronger
closure mechanism, while the old answer itself is absent from the dependency
closure of the new answer.

\subsection{The self-indexed tower and final family}

\begin{lstlisting}
def tower14 (n : Nat) : Transformer :=
  diagonalOperatorTower schedule14 fdFamilyTransformer (family13 n)

def F14 (n : Nat) : Nat :=
  tower14 n (fun _ => family13) n n
\end{lstlisting}

Three occurrences of $n$ serve different roles.

\begin{enumerate}[leftmargin=*]
\item $\mathsf{family13}(n)$ chooses the tower rank.
\item The output coordinate is $n$; every progressive transformer dominates
  all seed coordinates $i\le n$ there.
\item The numerical input is $n$, setting the syntax/fuel bound used whenever
  the dynamic-kind reflector is reached.
\end{enumerate}

The seed family is constant in its first coordinate:
$\mathrm{seed}(i,m)=\mathsf{family13}(m)$.  Choosing visible coordinate $i=0$
therefore exposes exactly $\mathsf{family13}(n)$ in the progressivity theorem.

\section{Formal comparison with Contender 13}

\subsection{Pointwise family comparison}

\begin{theorem}[\lean{family13\_lt\_F14}]
For every $n$,
\[
\mathsf{family13}(n)<F_{14}(n).
\]
\end{theorem}

\begin{proof}
The selected tower rank is progressive.  Apply progressivity to the constant
seed family, visible coordinate $i=0$, output coordinate $o=n$, and numerical
input $n$.  Since $0\le n$, the left side is
$\mathrm{seed}(0,n)=\mathsf{family13}(n)$, while the right side unfolds to
$F_{14}(n)$.
\end{proof}

The proof does not need to unfold the astronomical rank
$\mathsf{family13}(n)$.  It uses the theorem that every finite rank is
progressive, instantiated symbolically at that rank.

\subsection{Input growth}

The repository proves $n<\operatorname{reflectFd}(f,n)$ for every oracle $f$.
Therefore
\[
n<F_{13}(n)<F_{13}(F_{13}(n))=\mathsf{family13}(n)<F_{14}(n).
\]
This yields:

\begin{theorem}[\lean{input\_lt\_F14}]
For every $n$, $n<F_{14}(n)$.
\end{theorem}

\subsection{Domination of the reflected programs}

The rank-zero transformer is already the family-aware dynamic-kind reflector.
Since $0<\mathsf{family13}(n)$, the arbitrary-rank theorem places that entire
rank below the selected final rank.  Combining this path with
\lean{fvalue\_lt\_reflect14} gives the following concrete language-level
certificate.

\begin{theorem}[\lean{fvalue\_lt\_F14}]
For every $t\in\mathrm{CTm}_{\le n}$,
\[
\mathrm{value}_{\Fd}
  (\mathrm{pack}(i,m\mapsto\mathsf{family13}(m)),n,t)<F_{14}(n).
\]
\end{theorem}

Thus the comparison is not based only on the seed value.  Every bounded
program in the incumbent's strongest certified object language, evaluated
against the packed predecessor family, is strictly below the new family at the
same input.  The theorem is deliberately limited to the canonical rank-zero
reflection node; domination of programs at every internal node is left as a
stronger path-indexing problem.

\subsection{The required closed comparison}

Define
\[
\C{14}=F_{14}(F_{14}(42)).
\]
Then
\begin{align*}
\C{13}
 &=\mathsf{family13}(42)\\
 &<F_{14}(42) &&\text{by pointwise family comparison}\\
 &<F_{14}(F_{14}(42)) &&\text{by input growth}\\
 &=\C{14}.
\end{align*}

The Lean theorem is:

\begin{lstlisting}
theorem contender_13_lt_contender_14 : contender_13 < contender_14 := by
  simpa [contender_13, contender_14, family13] using
    (Nat.lt_trans (family13_lt_F14 42) (input_lt_F14 (F14 42)))
\end{lstlisting}

The proof has no numerical evaluation.  Lean only checks the definitions,
structural recursion, and inequalities.

\section{Why this is a significant rather than cosmetic increase}

A strict inequality alone can be cheap.  The present construction contains
several independently meaningful amplifiers.

\begin{center}
\small
\begin{tabular}{>{\raggedright\arraybackslash}p{0.19\textwidth}
                >{\raggedright\arraybackslash}p{0.32\textwidth}
                >{\raggedright\arraybackslash}p{0.37\textwidth}}
\toprule
Layer & Contender 13 & Contender 14 \\
\midrule
Object language & dynamic-kind System $F_\omega$ & same certified language \\
Oracle interface & one unary family from Contender 12 & packed countable family at every reflection node \\
Visible domination & current unary oracle value & every coordinate $fs(i,n)$ with $i\le n$ \\
Reflection geometry & one direct reflection & recursive operator stages \\
Schedules & none in final wrapper & rank-dependent, countably diagonal schedules above \lean{tree12} \\
Meta-rank & one reflection & $\mathsf{family13}(n)$ operator-stage levels \\
Formal margin & $\C{12}<\C{13}$ & all lower ranks are pointwise below all higher ranks; bounded dynamic-kind programs are below $F_{14}$; and $\C{13}<\C{14}$ \\
\bottomrule
\end{tabular}
\end{center}

The rank theorems are particularly important.  At input 42, the height is
exactly $\C{13}$.  A direct theorem exhibits every successor increase through
a concrete branch, and a second theorem closes those steps transitively to
compare arbitrary ranks.  Together with the program-domination theorem, the
construction contains an internally certified chain of strict upgrades indexed
by the predecessor's full construction.  No attempt is made to compress this into a speculative
single ordinal label.

\subsection{Why no exact fast-growing calibration is claimed}

The repository's later contenders combine impredicative normalization,
higher-kinded operators, a predicative dynamic universe, oracle reflection,
and countably branching meta-schedules.  Assigning a proof-theoretic ordinal to
that entire combination would require a separate ordinal analysis, including a
precise correspondence between the encoded object languages, the semantic
normalization arguments, and the external operator-stage recursion.  The
present paper therefore distinguishes three levels of assertion:

\begin{enumerate}[leftmargin=*]
\item \textbf{Kernel-level intended theorem:} the supplied Lean source states
  the exact strict inequalities described above.
\item \textbf{Structural mathematical claim:} the new number iterates a
  progressive dynamic-family reflector through a predecessor-indexed tower,
  with strict successor embeddings.
\item \textbf{Ordinal heuristic:} deliberately left open pending a dedicated
  calibration theorem.
\end{enumerate}

This is stronger scholarship than assigning an impressive but unsupported
ordinal expression.

\section{Computability and consistency discipline}

\subsection{Totality of the new definitions}

Every new recursion is visibly structural.

\begin{itemize}[leftmargin=*]
\item \lean{reflect14} is finite arithmetic plus finite list maxima and the
  existing total \lean{Fd.fvalue}.
\item \lean{schedule14} constructs one W-tree node and, on each branch, applies
  existing total syntax constructors and conversions.
\item \lean{diagonalOperatorTower} recurses structurally on a natural rank.
\item \lean{tower14}, \lean{F14}, and \lean{contender\_14} are ordinary
  transparent definitions.
\end{itemize}

Evaluation is fantastically infeasible but computable in principle.  This is
exactly the contest's intended distinction.

\subsection{No new normalization theorem is needed}

A tempting route beyond Contender~13 would be an impredicative universe that
contains itself, or kinds quantifying over kinds without a stratification.  Such
systems risk Girard-Hurkens-style inconsistency and would require a new strong
normalization proof before their program maxima could be used constructively.
Contender~14 avoids that risk.  It reuses the already established
\lean{Fd.sn\_of\_typed} theorem and changes only the oracle interface and the
external reflection geometry.

\subsection{Expected axiom closure}

The new arithmetic and transformer proofs themselves use ordinary constructive
reasoning.  Through the imported normalization development, the repository
expects the permitted axioms \lean{propext} and \lean{Quot.sound}.  The supplied
audit additions request \lean{\#print axioms} for the new number and principal
theorems and reject any forbidden predecessor constants in the definition
closure.

\section{No-laziness analysis}

The rule against laziness is not merely textual.  One could write a trivial
successor of the old closed number and obtain a formal inequality without
mathematical progress.  Contender~14 differs in four ways.

\begin{enumerate}[leftmargin=*]
\item The executable source of \lean{contender\_14} contains no reference to
  \lean{contender\_13} or any earlier closed contender constant.
\item The previous construction appears as a function
  \lean{family13 : Nat -> Nat}, not as a single numeral.
\item That family is used as the oracle seed of a stronger reflected language
  and as the rank of an operator-stage tower.
\item The strict successor-rank theorem is independent of the old contender:
  it applies to every seed family and every rank.
\end{enumerate}

The theorem \lean{tower14\_rank\_at\_42} is intentionally stated only in the
proof module:

\begin{lstlisting}
theorem tower14_rank_at_42 :
    tower14 42 =
      diagonalOperatorTower schedule14 fdFamilyTransformer contender_13 := by
  rfl
\end{lstlisting}

It documents the scale of the rank while keeping the closed predecessor
constant out of the new definition's dependency closure.

\section{Alternatives considered and rejected}

\subsection{A single second application of \lean{F13}}

Defining $F_{13}^{(k)}(42)$ for a small fixed $k$ would be easy and formally
larger, but it would not introduce a new mechanism.  It would also sit close to
the no-laziness boundary.  The present proposal instead changes the oracle
arity, reflection geometry, schedules, and meta-rank.

\subsection{One operator stage over the dynamic reflector}

A direct analogue of Contender~12 using the stronger dynamic reflector is a
reasonable intermediate candidate.  It is, however, dominated in design by
the diagonal tower: rank zero recovers the base reflector, rank one already
contains one operator stage, and every further rank is formally strictly
larger.

\subsection{Assuming oracle monotonicity}

This would simplify many comparisons but would be mathematically wrong.  The
explicit prefix maximum and canonical leaf branch were included precisely to
make every inequality direct.

\subsection{A self-containing dynamic universe}

This might appear to offer a more dramatic language-strength jump, but without
a stratified model and normalization proof it would jeopardize the constructive
basis of the program maximum.  A serious future contender should first build a
ranked universe hierarchy and prove its normalization, rather than silently
crossing into an inconsistent calculus.

\subsection{Using the predecessor constant as the tower height}

The definition could have used \lean{contender\_13} directly.  That would fail
the repository's dependency audit and obscure whether the prior construction
was genuinely integrated.  The chosen \lean{family13 n} formulation exposes
the underlying function uniformly at every input and supports the pointwise
theorem \lean{family13\_lt\_F14}.

\section{Integration and validation protocol}

The archive contains a unified patch and an overlay.  The intended integration
steps are:

\begin{enumerate}[leftmargin=*]
\item add \lean{BusyLean/Contender14.lean};
\item add \lean{BusyLean/Comparison14.lean};
\item import both from \lean{BusyLean.lean};
\item update \lean{scripts/check.py} so Contender~13 is the unbudgeted previous
  phase and Contender~14 is the current definition phase;
\item extend \lean{test/Audit.lean} with axiom prints, closure exclusions, and
  theorem checks;
\item run the static overlay validator;
\item run a clean \lean{python3 scripts/check.py} under Lean~4.34.0;
\item run \lean{lake env lean test/Audit.lean} independently and retain the
  full output.
\end{enumerate}

\subsection{What has been validated in this archive}

The following checks were executed locally:

\begin{itemize}[leftmargin=*]
\item all required source and theorem names are present;
\item no forbidden construct occurs in executable source text;
\item the new definition text does not mention any closed contender constant
  from Contender~5 through Contender~13;
\item the final definition and predecessor-family equations match the intended
  construction;
\item the unified patch is syntactically valid and applies to files with the
  recorded baseline contexts;
\item the LaTeX article compiles, and the resulting PDF is rendered and
  visually inspected;
\item all distributed files receive SHA-256 checksums.
\end{itemize}

\subsection{What has not been validated here}

The environment used to prepare this archive had neither a Lean~4.34.0 binary
nor a materialized checkout of BusyLean.  Consequently the following claims
are not made:

\begin{itemize}[leftmargin=*]
\item that the Lean elaborator accepted every source line;
\item that the kernel checked every theorem;
\item that \lean{\#print axioms} returned only the permitted axioms;
\item that the definition and proof phases met the 15-second and 60-second
  budgets.
\end{itemize}

These are integration gates, not optional polish.  The source was written to
mirror existing theorem signatures and proof patterns closely, but the archive
labels its status honestly as a serious, review-ready proposal rather than a
falsely certified build.

\section{Further research questions}

\subsection{Formal questions directly suggested by Contender 14}

\begin{enumerate}[leftmargin=*]
\item \textbf{Internal-node program domination.}  The supplied
  \lean{fvalue\_lt\_F14} theorem handles the canonical rank-zero reflection
  node.  Index arbitrary reflection nodes inside nested \lean{operatorStage}
  schedules and prove that every program value at every such node is below the
  root value $F_{14}(n)$.

\item \textbf{Quantitative rank-path certificates.}  Refine
  \lean{diagonalOperatorTower\_rank\_lt} so its proof returns an explicit
  branch path (or a list of local inequalities) witnessing how rank $a$ embeds
  into rank $b$.  This would make the large-margin certificate computationally
  inspectable rather than merely propositional.

\item \textbf{Schedule cofinality.}  Define an explicit embedding relation on
  W-tree schedules and prove that the positive branches of \lean{schedule14 k}
  are cofinal in the repository's finite \lean{Autonomous.hyper} hierarchy
  above \lean{tree12}.

\item \textbf{Definition-closure audit theorem.}  Turn the current metaprogram
  audit into a reusable theorem or verified command proving that construction
  families, but no closed predecessor constants, occur in each contender's
  dependency closure.

\item \textbf{Compilation-cost control.}  Measure whether the symbolic rank
  \lean{family13 n} stays opaque enough during elaboration to meet the 15-second
  definition budget.  If not, reorganize declarations to prevent accidental
  reduction while preserving transparency and executability.
\end{enumerate}

\subsection{Stronger language directions}

\begin{enumerate}[leftmargin=*,resume]
\item \textbf{Ranked dynamic universes.}  Replace one kind \lean{dyn} by a
  predicative hierarchy $\mathrm{dyn}_0,\mathrm{dyn}_1,\ldots$, where level
  $r+1$ packages operators over kinds formed using lower levels.  Prove strong
  normalization uniformly in $r$, add a sound checker, and reflect over a
  runtime-selected finite rank.

\item \textbf{Tree-valued object language.}  Extend the strongest normalized
  language so certified programs can return W-trees to Lean, not only Church
  numerals.  This could replace externally supplied schedules by schedules
  generated by the strongest language itself, echoing the self-indexing that
  made Contender~6 effective.

\item \textbf{Universe-polymorphic family coding.}  Internalize the packed
  family interface at the type level, so object-language programs can quantify
  over family representations and certified decoders rather than relying on a
  fixed 2-adic pairing.

\item \textbf{Proof-relevant normalization bounds.}  Extract a certified bound
  on normalization height from the reducibility proof and reflect over those
  bounds directly.  This may expose information currently erased by the normal
  form evaluator.
\end{enumerate}

\subsection{Calibration and metamathematics}

\begin{enumerate}[leftmargin=*,resume]
\item \textbf{Ordinal analysis of dynamic-kind $F_\omega$.}  Identify a theory
  whose provably total numerical functions match the accepted object language,
  including two-level kind variables and the predicative dynamic universe.

\item \textbf{External reflection versus internal provability.}  Determine
  whether the family reflector and operator-stage tower correspond to iterated
  uniform reflection principles over that theory, and under what coding
  assumptions.

\item \textbf{Conservative lower bounds.}  Embed standard fast-growing
  hierarchies or iterated inductive-definition systems into the object language
  to obtain rigorous lower bounds before attempting an exact calibration.

\item \textbf{Upper bounds from semantic ranks.}  Assign well-founded ranks to
  the reducibility-candidate semantics and the external W-tree recursion, then
  prove a compositional upper bound for the complete construction.
\end{enumerate}

\section{Conclusion}

Contender~14 is designed as the natural synthesis of the repository's two most
powerful recent ideas: Contender~13's predicative universe of operators and
Contender~12's operator-level reflection.  The new family reflector exposes a
countable oracle family to the strongest certified object language while
retaining direct visible-coordinate bounds.  A rank-dependent schedule then
iterates the operator closure, and the predecessor family itself selects the
number of ranks.  The explicit leaf branch provides a reusable strict-rank
witness; its finite transitive closure compares arbitrary ranks, and the
rank-zero path lifts bounded dynamic-kind program values below $F_{14}$.
Consequently the large-margin claim rests on structural inequalities rather
than oracle monotonicity or speculative ordinal slogans.

The central equations are compact:
\[
\mathsf{family13}(n)=F_{13}(F_{13}(n)),\qquad
F_{14}(n)=T_{\mathsf{family13}(n)}(\mathsf{family13})(n,n),\qquad
\C{14}=F_{14}(F_{14}(42)),
\]
and the intended formal conclusion is
\[
\boxed{\C{13}<\C{14}}.
\]
The accompanying overlay, patch, audit additions, validation script, and full
source listings make the proposal ready for kernel checking and repository
review.

\appendix

\section{Theorem dependency map}

\small
\begin{longtable}{>{\raggedright\arraybackslash}p{0.36\textwidth}
                  >{\raggedright\arraybackslash}p{0.56\textwidth}}
\toprule
Declaration & Role and principal dependencies\\
\midrule
\endfirsthead
\toprule
Declaration & Role and principal dependencies\\
\midrule
\endhead
\path{lt_reflect14} & Arithmetic fact for the outer $n+1+\max$ wrapper.\\
\path{oracle_lt_reflect14} & Uses \path{W.le_prefixMax} and \path{lt_reflect14}.\\
\path{fvalue_lt_reflect14} & Uses \path{W.le_maximum} and \path{lt_reflect14}.\\
\path{fdFamilyTransformer_progressive} & Immediate from visible-family domination.\\
\path{diagonalOperatorTower_progressive} & Natural-number induction plus the existing \path{operatorStage_progressive}.\\
\path{transformer_lt_operatorStage_schedule14} & Chains \path{transformer_lt_lift_at} with \path{operatorStage_child_lt} through branch zero.\\
\path{diagonalOperatorTower_step_lt} & Instantiates the preceding theorem at rank $k$ using tower progressivity.\\
\path{diagonalOperatorTower_rank_lt} & Inductively closes strict successor steps to compare arbitrary finite ranks.\\
\path{family13_lt_F14} & Applies tower progressivity to the constant predecessor seed at visible coordinate zero.\\
\path{input_lt_F13} & Reuses \path{Fd.input_lt_reflectFd}.\\
\path{input_lt_family13} & Two transitive uses of \path{input_lt_F13}.\\
\path{input_lt_F14} & Chains input growth of the predecessor family with \path{family13_lt_F14}.\\
\path{lower_rank_lt_F14} & Instantiates the arbitrary-rank theorem below the self-indexed final rank.\\
\path{fvalue_lt_F14} & Composes base-program domination with the rank-zero-to-final path.\\
\path{tower14_rank_at_42} & Definitional equality showing the exact meta-rank.\\
\path{contender_13_lt_contender_14} & Pointwise predecessor domination followed by input growth.\\
\bottomrule
\end{longtable}

\section{Lean definition module}

\lstinputlisting{../repo_overlay/BusyLean/Contender14.lean}

\section{Lean comparison module}

\lstinputlisting{../repo_overlay/BusyLean/Comparison14.lean}

\end{document}
